\chapter{Topología euclídea de $\R^n$}
De nuevo, asumimos que $V$ es un espacio vectorial sobre $\R$.
\begin{definition}[Punto interior]
  Llmamos a $x\in A\subseteq V$ un punto interior si existe $\epsilon >0$ tal que $B_\epsilon (x)\subseteq A$.
\end{definition}
\begin{definition}[Interior de un conjunto]
  Dado $A\subseteq V$, su interior es \[
      \Int A = \{x\in A : x \text{ es un punto interior de } A\} 
  \] 
\end{definition}
\begin{definition}[Conjunto abierto]
  Llamamos a $A\subseteq V$ abierto si $\Int A=A$.
\end{definition}
\begin{lemma}
  Toda bola es un conjunto abierto.
\end{lemma}
\begin{proof}
  Sea $x \in B_\epsilon (a)$ y $\epsilon '= \epsilon -\norm{x-a}>0$; considera la bola $B_{\epsilon'}(x)$, si $y\in B_{\epsilon'}(x)$, entonces $\norm{y-x}<\epsilon - \norm{x-a}$ y suando la desigualdad triangular \[
      \norm{y-a}\leq \norm{y-x}+\norm{x-a}< \epsilon ,
    \] lo cual implica $y\in B_\epsilon (a)$ y que para todo $x\in B_\epsilon (a)$, $B_{\epsilon'} (x) \subseteq B_\epsilon (a)$. Esto es, por definición, $\Int B_\epsilon (a)= B_\epsilon (a)$.
\end{proof}
\begin{proposition}[Propiedades de conjuntos abiertos]
  \begin{enumerate}
    \item La unión (finita o infinita) de una familia de conjuntos abiertos es un conjunto abierto.
    \item La intersección de una familia \emph{finita} de conjuntos abierto es un conjunto abierto.
    \item Si $A \subseteq B$, entonces $\Int A \subseteq \Int B$.
  \end{enumerate}
\end{proposition}
\begin{lemma}[El interior es el subconjunto abierto más grande]
  \label{lemma:conjunto_abierto_mas_grande_es_Int}
 Sea $A\subseteq V$, el conjunto abierto más grande contenido en $A$ es $\Int A$; esto es, si $C\subseteq A$ y $C$ es abierto, entonces $C \subseteq \Int A$.
 \end{lemma}
\begin{proof}
  Sea $x\in \Int C$, entonces para algún $\epsilon >0$, $B_\epsilon (x)\subset C$; como $C\subseteq A$, $x\in \Int A$; finalmente como $C$ es abierto $C=\Int C\subset \Int A$ como se quería demostrar.
\end{proof}
\begin{lemma}
    Sean $A,B\subseteq V$, entonces $\Int(A \cap B)= \Int A \cap \Int B$ 
 \end{lemma}
 \begin{proof}
   Como $\Int A, \Int B$ son conjuntos abiertos, su intersección $\Int A \cap \Int B$ lo es; como $\Int A \subseteq A$ y $\Int B \subseteq B$, entonces $\Int A \cap \Int B \subseteq A \cap B$. Luego, por \fullref{lemma:conjunto_abierto_mas_grande_es_Int}, sabemos que $\Int A \cap \Int B \subseteq \Int(A\cap B)$. La otra dirección de la inclusión es trivial.
 \end{proof}
 \begin{definition}
   Si $A\subseteq V$, denotamos su complementario rrespecto a $V$ por$A^c=V\setminus A$.
 \end{definition}
 \begin{definition}[Conjunto cerrado]
   Llamamos a $A\subseteq V$ cerrado si $A^c$ es abierto.
 \end{definition}
